\documentclass[12pt,a4paper,oneside]{article}
\usepackage[brazil]{babel}
\usepackage[utf8]{inputenc}
\usepackage{olymp}

\setcounter{problem}{4}

\begin{document}

\begin{problem}{MDC}{mdc}{2}
 
O \textbf{MDC}, Máximo Divisor Comum, entre dois números inteiros $A$ e $B$ é o maior número inteiro positivo que é divisor de $A$ e de $B$, ou seja, que tanto $A$ quanto $B$ são múltiplos dele.

O MDC é útil, por exemplo, para simplificação de frações.

Faça um programa para calcular MDC.

%\Notes
%
%Observações, se tiver.

\InputFile

A entrada é composta por dois valores inteiros, $A$ e $B$. Restrição: $1 \leq A, B \leq 10000$.


\OutputFile

Seu programa deve gerar apenas uma linha na saída, contendo o MDC de $A$ e $B$.


\Example

\begin{example}
\exmp{
20 30
}{
10
}%
\end{example}

\begin{example}
\exmp{
30 3
}{
3
}%
\end{example}

\begin{example}
\exmp{
99 77
}{
11
}%
\end{example}

\begin{example}
\exmp{
7 9
}{
1
}%
\end{example}

%Comentários sobre o exemplo, se necessário.

\end{problem}

\end{document}
